\documentclass[10pt,a4paper]{amsart}
\usepackage[margin=19mm]{geometry}
\usepackage{amsmath,amssymb,mathtools}
\usepackage[scr=rsfs,cal=euler]{mathalfa}
\usepackage{xfrac}
\usepackage[unicode,hidelinks,bookmarks=false]{hyperref}
\newcommand{\ie}{i.e.\ }
\newcommand{\cf}{cf.\ }
\newcommand{\resp}{respectively\ }
\newcommand{\Subsection}{\subsection}
\providecommand{\nobreakdash}{\nobreak\hbox{-}}
\setlength{\emergencystretch}{2em}
\multlinegap0pt
\usepackage{bm}
\usepackage{bbm}
\usepackage{dsfont}
\usepackage{xspace}
\usepackage{cancel}
\usepackage{mathtools}
\usepackage{amsthm}
\makeatletter
\@ifundefined{thm}{\newtheorem{thm}{Thm}}{}
\@ifundefined{prop}{\newtheorem{prop}[thm]{Proposition}}{}
\makeatother
\DeclareMathOperator{\Hst}{Hst}
\newcommand{\NN}{\mathbb{N}}
\title[Spherical Designs on $S^1$ of Finite Harmonic Strength]{Spherical Designs on $S^1$ of Finite Harmonic Strength}
\author{Ryutaro Misawa, Yusaku Nishimura}
\date{Source: arXiv:2505.06893v2 (CC BY 4.0)}
\begin{document}
\maketitle

\noindent\emph{Source and licence.} Spherical Designs on $S^1$ of Finite Harmonic Strength. Version arXiv:2505.06893v2 (CC BY 4.0), distributed under the Creative Commons Attribution 4.0 International licence, as recorded in the arXiv licence metadata for this exact version and shown on the abstract page https://arxiv.org/abs/2505.06893v2. Full rights notice: licenses/arxiv-2505.06893-CC-BY-4.0.txt.\par\smallskip

\noindent\textit{Editorial scope.} Counted unit: the Prop whose source label is \texttt{prop:23} in the article (source0.tex lines 510--514, source label \texttt{prop:23}), delivered together with the article's own complete original proof of it (source0.tex lines 516--599, one \texttt{proof} environment of 84 lines). No mathematics is written, repaired or re-derived: statement and proof are the article's own text line for line. Required context retained uncounted: prop:NT\_ge\_5 (lines 473--478). Every definition, display and label of those lines is retained unchanged. Omitted from this package and not redistributed: 1--472, 479--509, 515--515, 600--1000. Nothing inside a retained excerpt is omitted or altered: every retained body is delivered exactly as the article's lines are written, so no omission receipt of any kind accompanies this package. Citations: the retained excerpts carry no citation command at all, so no bibliography entry of the article is transcribed and no reference is redistributed. Reference adaptation: none was needed and none was made; every reference inside the delivered unit resolves inside it, so no unresolved reference or citation is produced. Printed slips kept literal: no printed word, symbol, hypothesis or number of the counted statement or of its proof was altered. Layout and class: the article's own class is \texttt{amsart} with packages including amsmath, amsthm, amssymb, mathtools, amscd, color, comment, cleveref, tikz, xy, enumerate, url, setspace; the wrapper is the lane's pinned \texttt{amsart} editorial wrapper at 10pt on A4, which restates the theorem-like environments the retained text needs and adds the article's own macro definitions that the retained text needs (\texttt{\textbackslash Hst}, \texttt{\textbackslash NN}), with extra wrapper packages (bm, bbm, dsfont, xspace, cancel, mathtools). The delivered numbering is the unit's own. Assets: no retained span contains a figure environment, a graphics-inclusion command, a PDF-page inclusion, a table or a TikZ picture; no image file is copied into this package, the package declares no asset dependency and no image file is redistributed with it. Licence: version arXiv:2505.06893v2 of this article is distributed under the Creative Commons Attribution 4.0 International licence, as recorded in the arXiv licence metadata for this exact version and shown on the abstract page https://arxiv.org/abs/2505.06893v2. A full rights notice is delivered at licenses/arxiv-2505.06893-CC-BY-4.0.txt. Characters: every retained excerpt is pure ASCII, so the delivered engine, XeLaTeX with its default Latin Modern text font, reports no missing character.
\begin{prop}\label{prop:NT_ge_5}
Let \(T\subset\NN\) be finite and nonempty. Then
\[
N(T,2)\ge5.
\]
\end{prop}

\begin{prop}\label{prop:23}
\[
N(\{2,3\},2)=5.
\]
\end{prop}

\begin{proof}
By Proposition~\ref{prop:NT_ge_5}, it is enough to construct a five-point
set.  Let
\[
X=\{1,e^{\pm ia},e^{\pm ib}\},
\qquad
x=\cos a,\quad y=\cos b.
\]
Then
\[
P_k(X)=1+2T_k(x)+2T_k(y),
\]
where \(T_k\) is the Chebyshev polynomial.  The equations
\(P_2(X)=P_3(X)=0\) are
\[
x^2+y^2=\frac34,
\qquad
8(x^3+y^3)-6(x+y)+1=0.
\]
Put \(s=x+y\) and \(p=xy\). Then
\[
p=\frac{s^2-\frac34}{2},
\qquad
(s-1)(2s+1)^2=0.
\]
Since \(P_1(X)=1+2s\), we take \(s=1\). Hence \(p=1/8\), and we may take
\[
x=\frac12-\frac{\sqrt2}{4},
\qquad
y=\frac12+\frac{\sqrt2}{4}.
\]
Then \(P_1(X)\neq0\) and \(P_2(X)=P_3(X)=0\).

It remains to show that there are no further vanishing moments. Put
\[
q_k:=P_k(X)-1.
\]
The four numbers \(e^{\pm ia},e^{\pm ib}\) are the roots of
\[
2z^4-4z^3+5z^2-4z+2.
\]
Therefore
\[
2q_{k+4}-4q_{k+3}+5q_{k+2}-4q_{k+1}+2q_k=0
\qquad(k\ge0).
\]
Moreover,
\[
q_0=4,\quad q_1=2,\quad q_2=q_3=-1,\quad q_4=\frac12,\quad q_5=-\frac12.
\]
For \(m\ge1\), define
\[
E_m:=2^{m-1}q_{2m},
\qquad
O_m:=2^{m-1}q_{2m+1}.
\]
Putting \(k=2m\) and \(k=2m+1\) in the recurrence gives
\begin{align*}
E_{m+2}-4O_{m+1}+5E_{m+1}-8O_m+4E_m&=0,\\
O_{m+2}-2E_{m+2}+5O_{m+1}-4E_{m+1}+4O_m&=0.
\end{align*}
Since
\[
E_1=O_1=-1,\qquad E_2=1,\qquad O_2=-1,
\]
these relations show by induction that \(E_m,O_m\in\mathbb Z\) for all
\(m\ge1\). Reducing them modulo \(2\), we obtain
\[
E_{m+2}\equiv E_{m+1},
\qquad
O_{m+2}\equiv O_{m+1}
\pmod 2.
\]
Hence \(E_m\) and \(O_m\) are odd for all \(m\ge2\).
Thus, for \(k\ge4\), the number \(q_k\) is not an integer. In particular
\(q_k\neq-1\). Since \(P_k(X)=0\) is equivalent to \(q_k=-1\), we get
\[
P_k(X)\neq0\qquad(k\ge4).
\]
Hence
\[
\Hst(X)=\{2,3\}.
\]
\end{proof}

\end{document}
